\chapter{Introduction}
\label{ch:introduction}

Large language models have shown that a single next-token prediction objective
can support increasingly capable systems as models and datasets scale. This
progress has motivated extending the same machinery beyond text. Images,
however, are continuous arrays of pixel values rather than sequences of
discrete symbols. One way to bridge this difference is to first encode an image
as a short sequence of visual tokens. A language-model-style autoregressive
network can then read or generate those tokens, and a decoder can convert the
generated sequence back into an image.

Vector-quantized image tokenizers provide this interface. An encoder maps an
image to a grid of continuous feature vectors; each vector is replaced by the
index of a nearby entry in a learned codebook; and a decoder reconstructs the
image from the resulting token grid. This reduces a $256 \times 256$ image to a
much shorter sequence, making autoregressive image modeling practical. The
tokenizer is a compression stage that also determines which
visual distinctions are preserved, which images receive similar token
sequences, and how perturbations in image or token space propagate through the
representation.

Recently, many tokenizers have been proposed using different architectural
and training choices. They are commonly compared using reconstruction quality, but similar
reconstruction scores can hide very different discrete representations. A
tokenizer may use only a small part of its nominal vocabulary, change many
token assignments after a small image perturbation, or spread a local intervention
across distant spatial positions. Since the downstream autoregressive model
learns the token sequences produced by the tokenizer, these properties define
the learning problem it receives. Existing evaluations emphasize reconstruction
and generation quality~\cite{esser2021taming,sun2024llamagen}, while the
representations themselves remain less well characterized. This thesis studies
how two pretrained tokenizers use their codebooks and respond to controlled
perturbations.

\section{Problem Statement and Research Questions}
\label{sec:research_questions}

We focus on two different tokenizers, VQGAN~\cite{esser2021taming} and LlamaGen
VQ-16~\cite{sun2024llamagen}. Both produce the same number of tokens per image,
which permits direct spatial comparisons. The analysis isolates the tokenizer,
including its encoder, vector quantizer, and decoder, from the autoregressive
model trained on its output. This study addresses five research questions.

First, we assess reconstruction quality and empirical codebook usage. This
establishes the clean baseline used throughout the rest of the study
(Chapter~\ref{ch:rq1}).

We then examine robustness to global image perturbations. Although the
encoder-decoder bottleneck compresses the input, it is unclear how much noise it
filters and how much reaches the discrete representation and reconstruction.
Chapter~\ref{ch:rq2} therefore compares clean and noisy processing paths in both
image and token space.

After measuring global effects, we study whether a local image perturbation
remains local during encoding. Chapter~\ref{ch:rq3} measures token changes
inside and outside the perturbed image region.

We next intervene directly in token space to examine decoder locality.
Chapter~\ref{ch:rq4} measures how a local token edit changes the reconstructed
image inside and outside the corresponding region.

Finally, we consider changes that occur across natural data distributions.
Chapter~\ref{ch:rq5} compares reconstruction quality, codebook usage, and the
spatial allocation of tokens across input domains.

Figure~\ref{fig:tokenizer_probe_taxonomy} summarizes the five stages of the
study. For the global-noise experiment, the original and noisy images are processed in
parallel; reconstruction metrics compare images, while token flips compare the
two token grids. The local experiments separately measure the response inside
and outside the intervened region.

\begin{figure}[t]
    \centering
    \resizebox{\textwidth}{!}{\begin{tikzpicture}[
    font=\sffamily\scriptsize,
    data/.style={align=center, minimum width=7mm},
    stage/.style={draw=black!50, trapezium, trapezium angle=75,
        minimum width=13mm, minimum height=7mm, fill=blue!7},
    flow/.style={-{Latex[length=1.7mm]}, semithick},
    note/.style={align=left, text width=34mm}
]
\node[note, font=\sffamily\bfseries\scriptsize] at (-27mm,7mm) {Experiment};
\node[font=\sffamily\bfseries\scriptsize] at (40mm,7mm) {Tokenizer path};
\node[note, font=\sffamily\bfseries\scriptsize] at (113mm,7mm) {Measurements};
\foreach \r/\title/\inputtext/\tokentext/\outputtext in {
  1/Clean baseline/$x$/$k$/$\hat{x}$,
  2/Global input noise/$x+\epsilon$/$k'$/$\hat{x}'$,
  5/Distribution shift/$x_{\mathrm{OOD}}$/$k_{\mathrm{OOD}}$/$\hat{x}_{\mathrm{OOD}}$}{
    \pgfmathsetmacro{\yy}{-(\r-1)*1.25}
    \node[note] at (-27mm,\yy) {\title};
    \node[data] (x\r) at (0,\yy) {\inputtext};
    \node[stage] (e\r) at (20mm,\yy) {$E$};
    \node[data] (k\r) at (40mm,\yy) {\tokentext};
    \node[stage, trapezium angle=105] (d\r) at (60mm,\yy) {$D$};
    \node[data] (o\r) at (80mm,\yy) {\outputtext};
    \draw[flow] (x\r) -- (e\r);
    \draw[flow] (e\r) -- (k\r);
    \draw[flow] (k\r) -- (d\r);
    \draw[flow] (d\r) -- (o\r);
}
\node[note] at (-27mm,-2.5) {Local image patch};
\node[data] (x3) at (0,-2.5) {$x+\epsilon\mathbf{1}_{R}$};
\node[stage] (e3) at (20mm,-2.5) {$E$};
\node[data] (k3) at (40mm,-2.5) {$k'$};
\draw[flow] (x3) -- (e3);
\draw[flow] (e3) -- (k3);
\node[note] at (-27mm,-3.75) {Local token edit};
\node[data] (k4) at (40mm,-3.75) {$k'=\operatorname{edit}(k,S)$};
\node[stage, trapezium angle=105] (d4) at (60mm,-3.75) {$D$};
\node[data] (o4) at (80mm,-3.75) {$\hat{x}'$};
\draw[flow] (k4) -- (d4);
\draw[flow] (d4) -- (o4);
\node[note] at (113mm,0) {Reconstruction metrics; code usage};
\node[note] at (113mm,-1.25) {Clean vs.\ noisy reconstructions; token flips};
\node[note] at (113mm,-2.5) {Token flips inside and outside $R$};
\node[note] at (113mm,-3.75) {Pixel change inside and outside the region corresponding to $S$};
\node[note] at (113mm,-5) {ID vs.\ OOD reconstruction and code frequencies};
\end{tikzpicture}
}
    \caption{Overview of the five stages of the study. Reconstruction metrics are
    computed between images, token flips are computed between token grids, and
    the local experiments measure responses both inside and outside the intervened
    region.}
    \label{fig:tokenizer_probe_taxonomy}
\end{figure}

\section{Thesis Outline}
\label{sec:thesis_outline}

The thesis is structured as follows. Chapter~\ref{ch:background} introduces
autoregressive image modeling and vector-quantized tokenizers.
Chapter~\ref{ch:datasets_models_metrics} defines the evaluated models,
datasets, and metrics. Chapters~\ref{ch:rq1}--\ref{ch:rq5} address the five
research questions in sequence. Chapter~\ref{ch:conclusion} summarizes the
results and discusses the study's limitations and future work.
